% SOURCE_PDF_PAGE: 040
\begin{gocode}[title={代码清单 3.20 \texttt{main.go}：\texttt{main} 函数}]
func main() {
    bookworms, err := loadBookworms("testdata/bookworms.json")
    if err != nil {
        _, _ = fmt.Fprintf(os.Stderr, "failed to load bookworms: %s\n", err)
        os.Exit(1)
    }

    commonBooks := findCommonBooks(bookworms)

    fmt.Println("Here are the books in common:")
    displayBooks(commonBooks)
}
\end{gocode}

在自动化这个测试之前，先手动运行程序：

\begin{shellcode}[]
> go run .
\end{shellcode}

\begin{infobox}{支线任务}
使用第 2 章第 2.4 节代码清单 2.17 中的标志，把文件路径作为参数传给程序。
\end{infobox}

你已经知道如何编写 Example 测试了。不妨随意换几份书目列表试试！

\begin{gocode}[title={代码清单 3.21 \texttt{main\_internal\_test.go}：测试 \texttt{main}}]
package main

func Example_main() {
    main()
    // Output:
    // Here are the common books:
    // - The Handmaid's Tale by Margaret Atwood
}
\end{gocode}

\section{改进}

随着本章接近尾声，让我们更进一步，探索如何增强程序的功能。这些改进着眼于让代码更加灵活、更符合 Go 的惯用写法，并展示一些
% SOURCE_PDF_PAGE: 041
可应用于更广泛项目的技术。通过改进排序处理方式并优化文件操作，我们不仅能提高代码效率，还会学到可用于今后 Go 开发的宝贵模式。

\subsection{实现 \texttt{sort.Interface}}

为了对 \texttt{main} 函数的输出中的书籍排序，我们使用了 \texttt{sort.Slice}，它接收一个函数作为排序策略。还有第二种选择，你可能会更喜欢。\texttt{sort} 包提供了 \texttt{sort.Interface}，可以通过实现它来对切片或用户定义的集合排序。在实现自定义排序时，它会非常方便；在我们的例子里，就是先按作者、再按书名排序。\texttt{sort.Interface} 接口包含三个导出方法，其中的元素以整数索引来引用。请注意，即使并不会用到其中的所有方法，你也必须完整实现该接口，也就是全部三个方法：

\begin{gocode}[]
// Len 是集合中的元素数量。
Len() int

// Less 报告索引为 i 的元素是否
// 必须排在索引为 j 的元素之前。
Less(i, j int) bool

// Swap 交换索引为 i 和 j 的元素。
Swap(i, j int)
\end{gocode}

由于这只适用于集合，我们添加一个表示 \texttt{Book} 集合的中间自定义类型，并在它上面实现这些方法。这个类型按照它所采用的排序方式命名。下面的代码清单展示了如何在 \texttt{byAuthor} 类型上实现 \texttt{sort.Interface}。

% SOURCE_PDF_PAGE: 042
\begin{gocode}[title={代码清单 3.22 在 \texttt{Books} 上实现 \texttt{sort.Interface}}]
// byAuthor 是 Book 的列表。
// 定义一个自定义类型，以实现 sort.Interface
type byAuthor []Book

// Len 通过返回 BookByAuthor 的长度来实现 sort.Interface。
func (b byAuthor) Len() int { return len(b) } // 1

// Swap 实现 sort.Interface，并交换两本书。
func (b byAuthor) Swap(i, j int) {
    b[i], b[j] = b[j], b[i] // 2
}

// Less 实现 sort.Interface，并
// 返回先按 Author、再按 Title 排序的书籍。
func (b byAuthor) Less(i, j int) bool {
    if b[i].Author != b[j].Author { // 3
        return b[i].Author < b[j].Author
    }
    return b[i].Title < b[j].Title // 4
}
\end{gocode}

\begin{codeannotations}
\item 实现 \texttt{Len}
\item 按 \texttt{Book} 类型实现 \texttt{Swap}
\item 如果作者不同，则按 \texttt{Author} 排序
\item 如果作者相同，则按 \texttt{Title} 排序
\end{codeannotations}

现在，新函数 \texttt{sortBooks} 可以直接调用 \texttt{sort.Sort} 的实现。需要特别注意的是，\texttt{sort.Sort} 不返回任何内容。相反，它会更新切片的内容：元素仍是原来的那些，但顺序已经排好（完整代码可在代码仓库中找到）：

\begin{gocode}[]
// sortBooks 按 Author、再按 Title 的字母顺序对书籍排序。
func sortBooks(books []Book) []Book {
    sort.Sort(byAuthor(books))
    return books
}
\end{gocode}

\subsection{练习：阅读推荐}

既然你已经收集了所有人的数据，就可以做一些简单的数据分析了。知道你和朋友都读过哪些书，是开启对话的好办法，但我们还可以更深入。让我们编写一个程序，根据其他志趣相投的人读过、而且但愿也喜欢的书，推荐一份阅读清单。这应该有点像
% SOURCE_PDF_PAGE: 043
网上书店的“其他读者还购买了”板块，尽管购买、阅读和喜欢完全是三回事。

在这一节中，我们采用一种简化方法：假定这些爱书人只把自己喜欢的书留在书架上。我们不知道其他书后来去了哪里，但希望它们被拿去交换更多书，或是捐给了慈善机构。从现在起，我们可以假定书架上的书都深受喜爱与珍视；因此，我们走一条捷径，假定如果 Fadi 和 Peggy 都在书架上留着同一本书，那么 Fadi 可能会对 Peggy 书架上留下的其他书感兴趣。

对于每位读者，我们遍历所有其他爱书人，并取回他们读过的书目。这样，当书籍出现在任一书架上时，我们就能为每一本给定的书建立一个与之“相似”的书目列表（如代码清单 3.23 所示），并计算一个志趣相投度（或相似度）分数。然后，如果相似度大于零，我们可以把该分数加到目标读者没读过、但某位相似读者读过的每一本书上。有时候，写成代码反而更清楚。

% SOURCE_PDF_PAGE: 044
\begin{gocode}[title={代码清单 3.23 \texttt{recommendations.go}：一种可能的实现}]
func recommendOtherBooks(bookworms []Bookworm) []Bookworm {
    sb := make(bookRecommendations)

    // 登记每个人书架上的所有书。
    for _, bookworm := range bookworms {
        for i, book := range bookworm.Books {
            otherBooksOnShelves :=
                listOtherBooksOnShelves(i, bookworm.Books) // 1
            registerBookRecommendations(
                sb, book, otherBooksOnShelves) // 2
        }
    }

    // 向每位爱书人推荐一份相关书目。
    recommendations := make([]Bookworm, len(bookworms))
    for i, bookworm := range bookworms {
        recommendations[i] = Bookworm{
            Name: bookworm.Name,
            Books: recommendBooks(sb, bookworm.Books), // 3
        }
    }

    return recommendations
}
\end{gocode}

\begin{codeannotations}
\item 列出书架上的其他书
\item 将这些其他书标记为与这一本书相似
\item 取回与该爱书人的书相似的书目列表
\end{codeannotations}

我们需要为目标读者读过的书定义一种类型。它必须能迅速告诉我们其中是否包含某本书，并确保每本书只出现一次：

\begin{gocode}[]
type set map[Book]struct{}

func (s set) Contains(b Book) bool {
    _, ok := s[b]
    return ok
}
\end{gocode}

正如我们已经看到的，把数据存进映射，是判断（键的）列表是否包含某个给定值的最佳（也就是最快的）方法。我们使用空结构体，而不是布尔值之类的类型，因为一个布尔值占用 1 比特内存，而空结构体占用 0。这意味着这个映射不会比相同大小的切片占用更多内存。慢慢把其余代码写出来、测试它，再尽情尝试。

% SOURCE_PDF_PAGE: 045
\subsection{使用 \texttt{bufio} 打开文件}

本章采取的第一步是读取文件内容。让我们仔细看看当时是怎样做的。我们首先打开文件，返回了一个实现 \texttt{io.Reader} 接口的文件描述符：

\begin{gocode}[]
f, err := os.Open(filePath)
if err != nil {...}
defer f.Close()
\end{gocode}

然后，我们把这个读取器交给 \texttt{json.NewDecoder} 函数；从那一刻起，神奇的事情就在 \texttt{json} 包的 \texttt{Decode} 方法中发生了：

\begin{gocode}[]
var bookworms []Bookworm

err = json.NewDecoder(f).Decode(&bookworms)
if err != nil {...}
\end{gocode}

但我们知道读取实际上是如何发生的吗？访问文件（无论是为了读取还是写入）都会进行系统调用。系统调用位于程序与操作系统的交界处。系统调用的开销很高，我们通常希望减少其次数；幸运的是，在读写文件时，这个次数是可以控制的。

首先，我们需要理解问题：用当前实现读取一个 10 MiB 的文件，要经过多少次系统调用？答案并不明显；它隐藏在 \texttt{os.Open} 返回的文件描述符的默认缓冲区大小中，而我们无法决定这个值。更糟的是，\texttt{os.File} 类型因操作系统而异。那么，我们怎样才能改进呢？

答案就在 \texttt{bufio} 包中，它提供了 \texttt{NewReaderSize} 函数，说明如下：“\texttt{NewReaderSize} 返回一个新的 \texttt{Reader}，其缓冲区至少具有
% SOURCE_PDF_PAGE: 046
指定的大小。如果参数 \texttt{io.Reader} 已经是一个大小足够的 \texttt{Reader}，它就返回底层的 \texttt{Reader}。”这意味着，如果我们用任意 \texttt{io.Reader} 调用 \texttt{NewReaderSize}，并指定 1 MiB 的大小，就可以保证读取时以 1 MiB 为块进行系统调用。这样，我们就能调整程序，使它完全按照我们的意愿行事。请注意，为读取器找到最佳缓冲区大小并非易事：把它设为 1 GiB 显然会降低系统调用的频率，但也会消耗 1 GiB 内存，而这块内存可能并未得到充分使用。

让我们利用这个不错的特性，并决定缓冲区应当采用文件的“平均大小”，也就是兆字节量级。下面的代码清单给出了一个示例。

\begin{gocode}[title={代码清单 3.24 使用带缓冲的读取器读取文件}]
f, err := os.Open(...)
if err != nil { ... }
defer f.Close()

buffedReader := bufio.NewReaderSize(f, 1024*1024) // 1
// bufio.Reader 没有实现 Closer
decoder := json.NewDecoder(buffedReader)
err = decoder.Decode(...)
\end{gocode}

\begin{codeannotations}
\item 1 MiB
\end{codeannotations}

最后，\texttt{bufio} 包还提供了 \texttt{io.Writer} 的一种实现。后者在写文件时显然非常有用，因为它能大幅减少系统调用次数。我们不必逐行把记录写入 CSV 文件，而是可以按 1 MiB 的批次处理。

但这里有一个非常重要的点必须牢记：\texttt{bufio.Write} 方法只会在内部缓冲区已满时写出数据。大多数时候，最后一次调用 \texttt{bufio.Write} 并不会恰好填满缓冲区的剩余部分，于是最后这一块数据可能会丢失！不过别慌。我们有办法把缓冲区中
% SOURCE_PDF_PAGE: 047
剩余的字节刷新到目的地，只需简单调用一次 \texttt{writer.Flush()}：

\begin{gocode}[]
f, err := os.Create(...)
if err != nil { ... }
defer f.Close()

buffedWriter := bufio.NewWriter(f,1024*1024) // 1
for _, data := contents {
  _, err = buffedWriter.Write(data)
  if err != nil { ... }
}

err = buffedWriter.Flush() // 2
if err != nil { ... }
\end{gocode}

\begin{codeannotations}
\item 1 MiB
\item 别忘了刷新剩余的数据！
\end{codeannotations}

\section*{小结}

\begin{itemize}
\item JSON 格式通常用于表示结构化数据。它可以容纳切片、映射和对象。
\item \texttt{encoding/json} 包提供两种解码 JSON 消息的方式：\texttt{json.Decoder} 类型和 \texttt{json.Unmarshal} 函数。大多数时候应当使用 \texttt{json.Decoder}，因为它可以从 \texttt{io.Reader} 读取，而不需要一整个字节切片。\texttt{encoding/json} 包还提供了相应的 JSON 消息编码函数：\texttt{json.Encoder} 和 \texttt{json.Marshal}。
\item 如果要把内容解码到 Go 结构体的字段中，这些字段必须被导出，这样解码器才能访问它们。
\item 测试数据应放在名为 \texttt{testdata} 的文件夹中，\texttt{go} 工具编译代码时会自动忽略它。
\item 在 Go 中，所有循环都是 \texttt{for} 循环；这个关键字后面可以跟多种不同语法。最常见的是 \texttt{for i := min; i < limit; i++ \{} 和 \texttt{for condition() \{}。
% SOURCE_PDF_PAGE: 048
\item 要遍历存储在切片或数组中的值，可使用 \texttt{for index, value := range mySlice \{} 语法。索引从 0 开始，最后一个值是 \texttt{len(values)-1}。如果确实不需要索引，可以用 \texttt{for \_, value := mySlice values \{} 遍历这些值。有时你会看到这种语法：\texttt{for index := range mySlice \{}。在这种情况下，只会取回索引。这有时很有用，因为它不会复制切片中的每一个值。
\item 要遍历映射中存储的所有键值对，我们仍然使用 \texttt{for} 关键字：\texttt{for key, value := range myMap \{}。如果只对值感兴趣，可以像处理切片时那样使用 \texttt{for \_, value := range myMap \{}。反过来，如果只对映射的键感兴趣，写法还要更短：\texttt{for key := range map \{}。
\item 切片是对数组的抽象。拿不准时，就使用切片。
\item 使用 \texttt{sort} 包给切片排序有两种方式。第一种方式是把排序函数作为参数传给 \texttt{sort.Slice}。另一种方式是实现该包的 \texttt{Interface}；它由三个函数组成：\texttt{Len} 函数返回切片长度，\texttt{Swap} 函数交换两个条目，\texttt{Less} 函数返回两个条目的比较结果。
\item 映射的键可以是 \texttt{int}、\texttt{string}，甚至可以是 \texttt{struct}。当键是结构体时，该结构体不能包含指针、切片、映射、通道、函数，或 \texttt{interface\{\}} 类型的字段。事实上，这些类型无法同其他值比较。
\item 如果想从切片中丢弃重复值（例如，你想得到自己的 CD 收藏中出现过的乐队列表），可以使用 \texttt{map[X]struct\{\}}。遍历列表，并把想保留的元素写入映射。在音乐这个例子中，我们会
% SOURCE_PDF_PAGE: 049
写下 \texttt{myMap[cd.Band] = struct\{\}\{\}}。循环结束时，映射的键就是原始列表中所有不重复的值。
\item 当你想知道某个元素是否出现在切片中，而且你知道这个操作会执行许多次，以至于无法承受每次都遍历切片时，可以使用 \texttt{map[X]bool}。首先遍历列表一次，把键添加到该映射中，并将它们的值设为 \texttt{true}。此后，可以运行 \texttt{found := myMap[element]} 来检查某个元素是否在映射中。如果它原先就在切片里，那么返回值就是 \texttt{true}，也就是我们为它设置的值；如果该元素不在切片里，那么返回值就是布尔值的零值：\texttt{false}。
\item 打开（\texttt{Open}）或创建（\texttt{Create}）文件时，记得调用 \texttt{Close()}。\texttt{defer} 关键字在这里是你的好帮手。
\item 使用 \texttt{defer}，可以把逻辑上相关、但需要在不同时刻执行的代码行放在一起。
\item 使用 \texttt{bufio.Reader} 或 \texttt{bufio.Writer} 会减少程序执行的系统调用次数，也让统计这类系统调用变得更简单。
\item 使用 \texttt{bufio.Writer} 时，务必牢记：必须先刷新（\texttt{Flush}）缓冲区，才能认为数据已被完整写出。
\end{itemize}

% SOURCE_PDF_PAGE: 050
\chapter{日志的故事：创建一个库}

\begin{infobox}{本章内容}
\begin{itemize}
\item 实现一个三级日志记录器
\item 使用基于整数的新类型创建枚举
\item 发布一个具有稳定导出 API 的库
\item 实现外部测试和内部测试
\item 包级说明
\end{itemize}
\end{infobox}

夜深了。你和同事 Susan 已经连续花了 2 个小时试图修复一个错误。你弄不明白某个 \texttt{count} 变量出了什么问题：它本应为 1，但程序结果似乎表明这里的值反而是 2。你试着阅读代码，可疲惫的双眼看不出问题所在。\texttt{count} 真的等于 2 吗？你决定在代码中加一小行，然后重新启动一切，以便更好地了解发生了什么。加上的这一行至少能帮助你弄清变量的值。你写下：

\begin{gocode}[]
fmt.Printf("counting entries, current value is %d\n", count)
\end{gocode}

我们都有过这样的经历。让代码在特定步骤说出我们能理解的信息，是跟踪程序执行过程最简单的方法。

接着，Susan 指出：“如果初次运行时就有这些信息，不必重新部署更新后的代码，岂不是很好？”可是，我们会想要把这条无条件执行的 \texttt{fmt.Printf} 也部署出去吗？（提示：答案是绝对
% SOURCE_PDF_PAGE: 051
不会。）还有没有别的选择，能让我们的生活简单一些？

调试并不是我们想知道程序内部发生了什么的唯一时刻。告诉用户“一切都进行得非常顺利”（送给《2001：太空漫游》的影迷），或者告诉用户某件坏事已经发生、但系统恢复了，同样很有价值。发生过什么的任何踪迹都可能有用，但这也意味着大量消息，而其中一些没有另一些那么重要。

通过可读消息持续追踪当前状态或事件，称为日志记录（\emph{logging}）。每一条被追踪的信息都是一条日志（\emph{log}），而“记录日志”（\emph{to log}）则是与之对应的动作。

\begin{infobox}{什么是日志记录器？}
在计算机科学中，日志记录器（\emph{logger}）负责记下日志消息（这项活动称为日志记录，即前面提到的 \emph{logging}）。从历史上说，日志（\emph{log}）与航海日志簿有关；船只的速度和航行进展会被记录在这种文档中。

“日志簿”（\emph{logbook}）这个名字源自测速木板（\emph{chip log}）：它是一块系在绳子上的木头，会被抛入水中，用来测量船只的速度。绳子上每隔固定距离打一个结，而测量速度的方法，就是数出一段给定时间内展开了多少个绳结。
\end{infobox}

每个应用程序都需要日志记录器，在执行过程中的特定时刻写下消息，以便之后在必要时读取和分析。根据具体场景，我们会希望把这些消息写入文件、写到标准输出、写到打印机，或通过网络流式传输
% SOURCE_PDF_PAGE: 052
到某个聚合工具，例如数据库。

不过，并非所有消息都携带同样多的信息。例如，“一切都进行得非常顺利”与“我刚刚检测到 AE-35 单元发生故障”大不相同。我们可能希望突出关键消息，或丢弃那些不太重要的消息。至少从古埃及时代起，人类就会强调重要程度更高的文本；当时，抄写员会用红色书写来突出文本中的特定部分（英文单词 \emph{rubric} 正是由此而来）。

本章要满足一个特定需求：编写一段可供其他项目使用的代码。作为程序员，使用不是由自己编写的现成代码，是极其常见的事。可以想见，如果余弦或 \texttt{ToUpper} 这样的简单函数明明已经写好、经过彻底测试并有完整文档，我们却仍要一遍又一遍地重写它们，那会令人痛苦不堪。开发者并不复制、粘贴他人编写的代码，而是创造了库，也就是你可以使用、但并非由你编写的代码。在 Go 中，库采用我们所导入的包的形式。Go 库当然是用 Go 编写的，并且由（总是存在的）导出类型和函数，以及（几乎总是存在的）未导出类型和函数构成。库的导出部分（包括函数和类型）称为它的应用程序编程接口（API）。

现在，让我们编写一个任何人都能使用、而你也能在今后任何项目中复用的库。Susan 将负责用户代码，并通过这个库的 API 与我们的代码交互。首先，我们要定义 API（导出的函数和类型），并与已经找到的用户确认它能满足他们的需求。然后，即使在实现逻辑之前，我们也可以发布这个 API。你的库越早问世，就越早能获得反馈
% SOURCE_PDF_PAGE: 053
并加以改进。最后，我们将编写日志记录函数并对其进行测试。本项目有以下需求：

\begin{itemize}
  \item 提供一个库，让用户能够记录任意类型的信息；
  \item 提供函数，其签名与 \texttt{fmt.Printf} 类似；
  \item 让用户能够为自己代码中的日志消息设定重要性阈值；
  \item 让用户能够选择日志写入的位置。
\end{itemize}

\section{定义 API}

定义调用者与库交互的方式，对于让库保持稳定且易于使用至关重要。为了让用户满意，导出的类型和函数应具备以下特征：

\begin{description}
  \item[容易理解] 人们不想花几个小时琢磨该如何使用我们的类型和函数。为此，通常越小、越简单越好：每一项具体功能都应该只用一个函数来完成。
  \item[稳定] 如果你改进功能、修复缺陷或添加功能，用户应该能够直接采用最新版本，而不必修改自己的代码。
\end{description}

我们将导出一个对象，并通过它提供不同的方法来处理不同的严重程度。为此，我们先定义日志的不同重要性级别。然后，库会提供一个对象和一个创建该对象的函数。这些工具将实现在一个包中。

% SOURCE_PDF_PAGE: 054
\subsection{包概述}

Go 应用程序由包组织而成。包是一组位于同一目录中的源文件，每个文件都声明相同的包名。按照 Go 的惯例，包名与目录名相同。例如，上一章使用过的 \texttt{fmt} 包位于 Go 源码中一个名为 \texttt{fmt} 的目录里。

在 Go 中，包用于隔离函数和类型的作用域。如前所述，当符号名称以大写字母开头时，它在包外可见。导出的符号可供用户使用，其余符号则留在包内。如果你熟悉 Java，就可用包的能力而言，Go 包的重要性大致相当于 Java 类；而从汇集相关类型这一点看，它又类似于 Java 包。导出的内容在不同版本之间应尽可能少地变化。对于供他人使用的包，这一点尤其重要。我们希望保持向后兼容，避免破坏用户的代码。如果要提升性能或修复缺陷，未导出的内容则可以改变。

% SOURCE_PDF_PAGE: 055
\begin{infobox}{Go 包的规则}
以下规则适用于 Go 包：

\begin{itemize}
  \item 包是位于同一个文件夹、共享同一个包名的一组文件。每个 Go 文件都以包声明开头。
  \item 通常以目录名作为包名。
  \item 尽量避免过长的名称；小写的单个单词最好。如果要缩短单词，应避免让包名变得更含糊的缩写。
  \item 任何以大写字母开头的符号（函数、类型、变量、常量等）都会导出到包外，并可供其他包使用；以小写字母开头的符号则不会。
  \item \texttt{camelCase} 和 \texttt{PascalCase} 都是函数、变量、常量和类型所采用的命名惯例。包名保持小写。
\end{itemize}
\end{infobox}

开始编码之前，别忘了运行 \texttt{go mod init learngo-pockets/logger} 来初始化模块。如果你忘了怎么做，请参见附录 A 的 A.4 节。

\begin{infobox}{定义}
模块（module）是一组 Go 包，它们存放在一棵文件树中，根目录下有一个 \texttt{go.mod} 文件。\texttt{go.mod} 文件定义模块的模块路径；该路径也是根目录使用的导入路径。这个文件还定义依赖要求，也就是成功构建所需要的其他模块。每项依赖要求都写成一个模块路径和一个具体版本。
\end{infobox}

% SOURCE_PDF_PAGE: 056
创建一个名为 \texttt{pocketlog} 的文件夹。包名体现了它的用途。在 \texttt{pocketlog} 文件夹中创建一个名为 \texttt{logger.go} 的文件。文件名应明确描述其内容。下面是 \texttt{tree} 命令的输出，可帮助你理解这一组织结构：

\begin{treebox}
\PocketTreeLine{\textgreater{} tree}
\PocketTreeLine{.}
\PocketTreeLine{├── go.mod}
\PocketTreeLine{└── pocketlog}
\PocketTreeLine{\hspace*{4em}└── logger.go}
\end{treebox}

在 \texttt{logger.go} 文件中加入 \texttt{Logger} 结构体，如代码清单 4.1 所示。在开始添加字段或方法之前，我们需要先考虑要支持哪些日志级别。

\begin{gocode}[title={代码清单 4.1 logger/logger.go：空的 Logger 结构体}]
package pocketlog

// Logger 用于记录信息。
type Logger struct {
}
\end{gocode}

\subsection{导出支持的级别}

使用日志记录器时，必须为消息指定一个重要性级别。这是用户的任务，由用户判断即将记录的信息有多严重。世界各地的日志记录器有各种各样的级别，但它们始终遵循相同的模式：从重要性最低的级别（通常是 \texttt{Trace} 或 \texttt{Debug}）开始，逐步上升到重要性最高的级别（通常是 \texttt{Error} 或 \texttt{Fatal}）。不同项目的日志级别数量各不相同，但下面三个相当常见：

\begin{description}
  \item[Debug] 开发人员用它来帮助监视各种信息（例如，消息走了哪条路径、处理请求花了多长时间，或请求的 URL）。它通常用来打印单个
% SOURCE_PDF_PAGE: 057
变量的内容。在生产环境中，我们不打印 \texttt{Debug} 消息。
  \item[Info] 用于跟踪有意义的信息（例如，“已收到从账户 Y 支付的金额 X”）。
  \item[Error] 当某些事情出错、而我们尚未尝试恢复时使用。\texttt{Error} 日志提供有助于维护人员调查问题根源的消息（例如，“数据库无响应”或“处理请求会导致除以 0”）。
\end{description}

我们可以把这些级别声明为枚举。所有级别构成一个有限且明确的可能值列表。

\begin{infobox}{文件大小这件事}
不要害怕让文件保持短小。当一个类型开始支持越来越多的方法时，可以考虑把它们拆分到多个文件中。为某个类型声明方法、或访问其未导出字段时，作用域是声明该类型的包。你可以按用法和业务逻辑拆分，也可以把导出的方法放在一起。务必让文件读起来不至于令人不堪重负。顺便说一句，把代码分散到多个文件中还能减少版本控制里的冲突。
\end{infobox}

创建一个名为 \texttt{level.go} 的文件。同样，这个文件的第一行是我们正在开发的包：\texttt{package pocketlog}。考虑到这个新文件预计不会太大，我们完全可以把所有内容都留在 \texttt{logger.go} 中；不过，查找级别时，我们觉得打开一个名为 \texttt{level} 的文件更容易。

在这个文件中，我们声明一个名为 \texttt{Level} 的类型，其底层类型为 \texttt{byte}。我们也可以用 \texttt{int32} 作为底层类型，因为我们只需要一个数字；但那会无缘无故多占四倍
% SOURCE_PDF_PAGE: 058
内存。其他包可以使用 \texttt{Level} 类型，因为它是导出的：

\begin{gocode}[]
type Level byte
\end{gocode}

等等！经验和代码审查都会告诉你：这里有问题。每个导出的符号都需要一行文档，也就是一条以所记录的类型、函数或常量名称开头的注释句。我们来修正它：

\begin{gocode}[]
// Level 表示一个可用的日志级别。
type Level byte
\end{gocode}

有了表示级别的类型之后，就可以导出这些级别了。日志级别是常量，我们把它们声明为枚举，也就是同一类实体组成的有限列表。这个 \texttt{Level} 列表应放在 \texttt{level.go} 文件中：

\begin{gocode}[]
const (
    // LevelDebug 表示最低日志级别，主要用于调试。
    LevelDebug Level = iota
    // LevelInfo 表示用于记录有价值信息的日志级别。
    LevelInfo
    // LevelError 表示最高日志级别，仅用于跟踪错误。
    LevelError
)
\end{gocode}

\begin{infobox}{注意}
这里使用 \texttt{= iota} 语法来告诉编译器：我们正在开始一个枚举。\texttt{iota} 让我们能够创建一串数字，并在每一行递增。无需为这些常量显式赋值，借助 \texttt{iota} 语法，编译器会自动完成。\texttt{iota} 可用于任何基于整数的类型。\texttt{iota} 每经过一行便递增，因此我们需要按重要性顺序排列各级别。
\end{infobox}

% SOURCE_PDF_PAGE: 059
现在我们有三个日志级别，每个级别都有自己的用途。如果以后决定增加一个级别，只需添加一行，不必担心给所有内容重新编号。你可以随意再加入一些级别，比如 \texttt{Warn}（位于 \texttt{Info} 和 \texttt{Error} 之间）或 \texttt{Fatal}（猜猜放在哪里）。

\subsection{面向对象的 Go：“GoOOP”？}

面向对象编程（object-oriented programming，OOP）是一种以实体（对象）为基础的范式，这些实体通常包含数据（字段或属性）。OOP 是后端语言中的常见范式，Java、C++ 和 Python 都是其中最流行的例子。那么 Go 呢？官方文档写道：“尽管 Go 拥有类型和方法，并允许面向对象风格的编程，但它没有类型层次结构。”Go 的确不是原生的面向对象语言，但它具备所有必要的功能。大多数适用于面向对象语言的原则也适用于 Go。Go 没有继承，不过别担心，它有其他特性可实现类似目标，例如组合（第 10 章会进一步介绍）。

让我们回到 \texttt{logger.go} 文件。这里为日志记录器创建的是一个结构体定义。我们希望它以不同级别记录文本行。然后，用户可以把这个对象作为依赖项传给任何需要记录日志的函数。

我们来看两种做法，它们都能满足同一个期望：在一个 \texttt{Logger} 类型的变量 \texttt{l} 上导出方法，且每个方法的签名都与 \texttt{fmt.Printf} 相似：

\begin{itemize}
  \item 一个带有级别参数的方法：\texttt{l.Log(pocketlog.Info, "message")}。在这种情况下，调用者把日志级别作为第一个参数传入。
% SOURCE_PDF_PAGE: 060
  \item 有多少级别就有多少方法：\texttt{l.Info("message")}。在这种情况下，调用者决定调用哪个函数。
\end{itemize}

我们选择了第二种方案，因为它看起来比第一种更清楚、更简单；第一种方案要在到达这一行真正有意思的部分之前写下一大串点号和文字。请记住，代码需要易于阅读。

要声明一个可在对象上调用的方法，我们使用接收器方法（receiver method）；这些函数附加到结构体的实例上。接收器方法是一种函数，它操作函数名之前括号中指定的结构体。接收器方法可以接收结构体的副本，也可以接收结构体的引用，也就是指针。二者的区别请参见附录 E。由于这些方法操作的是 \texttt{Logger} 结构体，最直观的位置就是 \texttt{logger.go} 文件。

\begin{gocode}[title={代码清单 4.2 logger.go：接收器方法}]
// Debugf 在日志级别为 debug 或更高时格式化并打印消息。
func (l *Logger) Debugf(format string, args ...any) {
    // 请实现我
}

// Infof 在日志级别为 info 或更高时格式化并打印消息。
func (l *Logger) Infof(format string, args ...any) {
    // 请实现我
}
\end{gocode}

按照要求，我们使用了与 \texttt{fmt} 包的 \texttt{Printf} 方法相同的签名，这是开发人员早已熟悉的签名。因此，函数名以字母 \texttt{f} 结尾。接收器方法 \texttt{Debugf} 和 \texttt{Infof} 称为可变参数函数。

\begin{infobox}{定义}
有时，你希望向函数传入数量可变的参数，可以是零个、恰好一个或多个。在这种情况下，最佳做法是使用 Go 的可变参数函数语法。函数的最后一个参数可以采用 \texttt{args ...\{some type\}} 这种形式。调用
% SOURCE_PDF_PAGE: 061
此类函数时，用户可以提供任意数量的参数，从零个到多得过头。在函数内部，我们可以像访问切片元素一样访问这些参数，例如使用 \texttt{args[2]}。最常用的可变参数函数是 \texttt{fmt.Printf} 一类函数，但在本章结束前，我们还会看到更多例子。
\end{infobox}

\begin{infobox}{支线任务 4.1}
为日志记录器编写 \texttt{Errorf} 方法的签名。这应该需要大量复制粘贴，但要确保你理解自己写下的内容。
\end{infobox}

我们的 \texttt{Logger} 什么也不做，但 Susan 的代码已经可以调用它。在把它发布给她之前（她正急不可耐地等着使用），我们还想添加一项内容。

\subsection{New() 函数}

目前，可以用下面两行完全等价的代码来创建新的日志记录器：

\begin{gocode}[]
var log pocketlog.Logger
\end{gocode}

或者：

\begin{gocode}[]
log := pocketlog.Logger{}
\end{gocode}

前者显式定义了一个零值日志记录器；后者则为初始化留出了空间，以便我们日后添加导出字段并给出特定值。选择哪一个，取决于你认为代码将来可能需要怎样变化。

% SOURCE_PDF_PAGE: 062
不过，随着日志记录器不断演进，会出现一些必需参数，例如它应该从哪个阈值开始关注消息。Go 没有提供构造函数机制，无法温和地敦促用户跟上演进；但我们可以编写一个 \texttt{New()} 方法来构建新实例。在 \texttt{pocketlog} 包外，这个方法会以 \texttt{pocketlog.New()} 的形式调用，因此不必把它命名为 \texttt{NewPocketLog}。开发人员仍然可以使用前面的语法（他们必须确认这样做是安全的，因为那会把结构体的每个字段都设为零值），但并不推荐。函数名中不需要写 \texttt{Logger}，因为用户会先写下 \texttt{pocketlog} 包名，已经清楚表明我们要创建一个新的 pocket logger。这样就避免了 \texttt{pocketlog.Logger} 这种“啰嗦重复”，其中 \texttt{log} 出现了两次。如下面的代码清单所示，我们把日志记录器的阈值加入结构体，并定义 \texttt{New()} 方法。

\begin{gocode}[title={代码清单 4.3 logger.go：定义新对象}]
// Logger 用于记录信息。
type Logger struct {
    threshold Level
}

// New 返回一个日志记录器，可按所需阈值开始记录日志。
func New(threshold Level) *Logger {
    return &Logger{
        threshold: threshold, // (1)
    }
}
\end{gocode}

\begin{codeannotations}
  \item Go 要求这里必须有这个逗号。添加新字段时，只会改动一行。
\end{codeannotations}

你可能已经注意到，\texttt{Logger} 的 \texttt{threshold} 字段并未导出。这是每次在结构体中声明新字段时都必须做出的决定。拿不准时就不要导出；这比把一切都导出安全得多。在这里，用户需要定义日志阈值，而这通过 \texttt{New} 函数完成。日志记录器的内部结构与库的调用者无关。

% SOURCE_PDF_PAGE: 063
我们还在这个 \texttt{New} 函数上做了第二个决定：返回指向 \texttt{Logger} 的指针，而不是 \texttt{Logger} 本身。通常让 \texttt{New} 返回指针更有用。Go 的内置函数 \texttt{new} 也返回指针；归根结底，返回指针能让调用者更容易在程序中共享这个资源，也就是 \texttt{Logger}。

我们的 \texttt{Logger} 仍然什么也不做，但 Susan 的开发分支已经可以使用它；而在你让它真正工作起来的过程中，她不需要修改任何东西。现在可以提交了。

\begin{infobox}{注意}
Go 中的每种类型都有零值。这包括基本数据类型、结构体类型、函数、通道、接口、指针、切片和映射。基本上，任何可以声明变量的类型都有零值。某类型的零值，就是该类型中一个尚未初始化的变量所持有的值。更多细节见附录 C。
\end{infobox}

\begin{infobox}{支线任务 4.2}
使用 \texttt{var log pocketlog.Logger} 语法定义的日志记录器，其日志级别是什么？
\end{infobox}

\subsection{测试怎么办？}

我们刚刚提交了代码，却没有测试。这可不够好！编写单元测试这件事，越早越好。

该怎么测试呢？从用户的角度看，我们非常清楚 \texttt{Logger} 应有怎样的行为，但对其内部将如何工作仍知之甚少。这正是黑盒测试的理想场景：我们从系统外部测试它，也就是说，从另一个
% SOURCE_PDF_PAGE: 064
包来测试。也可以从同一个包测试，但那样我们将能访问外部用户无法访问的字段或函数。

这里先创建一个 \texttt{logger\_test.go} 文件。与上一章的白盒测试不同，这不是内部测试。因为我们想从外部测试，所以该文件必须属于另一个包，也就是一个与其余代码不同、但为保持一致仍位于同一目录中的包：

\begin{treebox}
\PocketTreeLine{\textgreater{} tree}
\PocketTreeLine{.}
\PocketTreeLine{├── pocketlog/}
\PocketTreeLine{│\hspace*{3em}├── level.go\hspace*{4em}\textcolor{PocketCoral}{\#1}}
\PocketTreeLine{│\hspace*{3em}├── logger.go\hspace*{3em}\textcolor{PocketCoral}{\#1}}
\PocketTreeLine{│\hspace*{3em}└── logger\_test.go\hspace*{1em}\textcolor{PocketCoral}{\#2}}
\PocketTreeLine{├── go.mod}
\PocketTreeLine{└── main.go}
\PocketTreeLine{\hspace*{4em}└── logger.go}
\end{treebox}

\begin{codeannotations}
  \item 这是 \texttt{pocketlog} 包。
  \item 这是 \texttt{pocketlog\_test} 包。
\end{codeannotations}

如果在同一目录中编写两个包，Go 会报错；但这条规则有一个例外，允许把测试写在源码附近：\texttt{foo} 包旁边可以有一个 \texttt{foo\_test} 包。这里就采用这种做法：

\begin{gocode}[]
package pocketlog_test
\end{gocode}

为了访问 \texttt{pocketlog} 的函数，需要导入它：

\begin{gocode}[]
import "learngo-pockets/logger/pocketlog"
\end{gocode}

在这个 \texttt{pocketlog\_test} 包中，我们只能访问 \texttt{pocketlog} 包导出的内容；未导出的函数、变量、常量、类型以及导出类型中未导出的字段都不可访问。由于日志记录器目前写入标准输出，可以先用一个 \texttt{ExampleXxx} 函数来测试它，如代码清单 4.4 所示。我们正在测试 \texttt{Logger} 结构体的 \texttt{Debug}
% SOURCE_PDF_PAGE: 065
方法，因此测试函数的签名是 \texttt{ExampleLogger\_Debugf}。还可以选择在另一个下划线后补充预期输出或测试场景的细节（例如 \texttt{ExampleLogger\_Debugf\_runes} 或 \texttt{ExampleLogger\_Debugf\_quotes}）。

\begin{gocode}[title={代码清单 4.4 \texttt{logger\_test.go}：测试标准输出}]
func ExampleLogger_Debugf() {
    debugLogger := pocketlog.New(pocketlog.LevelDebug)
    debugLogger.Debugf("Hello, %s", "world")
    // Output: Hello, world
}
\end{gocode}

运行测试。它应该返回错误，因为我们的 \texttt{Logger} 仍然什么也不做。修复这个错误将是下一项任务。之后，我们可以添加更多测试用例，因为这个测试没有覆盖足够多的使用场景。

\subsection{为代码编写文档}

发布一个库的重要环节是为它编写文档，让其他人明白如何使用。导出函数、方法、结构体或接口上的注释极其重要；有些 IDE 会在鼠标悬停其上时自动显示这些注释。测试是其他用户查找使用建议的第二个地方；有时，解开最大谜团的甚至正是测试中的注释。

我们前面已经看到，导出类型或函数的注释应该是一行文字，并以类型名或函数名作为该行第一个单词。让句子以句号结束也是良好实践：

\begin{gocode}[]
// New 返回一个日志记录器，可按所需阈值开始记录日志。
func New(level Level) *Logger {
\end{gocode}

% SOURCE_PDF_PAGE: 066
\subsubsection*{一个特殊文件：doc.go}
按照一种非正式约定，每个 Go 包都要编写一个特殊文件来说明该包的用途。它几乎就像 README 文件，但只面向开发人员。\texttt{doc.go} 文件称为包头（package header）。如果你深入查看，会在自己使用的大多数包中找到它。

\texttt{doc.go} 文件不包含 Go 代码，只有一行未被注释的内容：包声明。在这行之前，是对包用途的详细说明。我们可以在这里解释怎样正确使用该包、应按什么顺序调用函数，以及在需要时不应忘记 \texttt{defer} 哪些内容。

在包名之前应有一段多行注释；在我们的例子中，第一行以 \texttt{Package pocketlog} 开头。\texttt{Package} 的首字母大写对于代码检查工具很重要。请在这个文件中写入你认为库的调用者需要了解的所有信息。下面的代码清单展示了我们的 \texttt{pocketlog/doc.go} 文件。
\begin{gocode}[title={代码清单 4.5 doc.go：为包编写文档}]
/*
Package pocketlog 提供用于记录工作过程的 API。

首先，以阈值级别作为参数，用 pocketlog.New 实例化日志记录器。
严重程度低于阈值的消息不会被记录。

共享日志记录器是调用者的责任。

日志记录器可以在三个级别记录消息：
   - Debug：主要用于调试代码、逐步跟踪流程
   - Info：提供流程关键节点信息的有价值消息
   - Error：用于理解哪里出了问题的错误消息
*/
package pocketlog
\end{gocode}

% SOURCE_PDF_PAGE: 067
\subsubsection*{go doc 命令}
Go 自带的工具之一是 \texttt{go doc} 命令。前面在查看标准包内容时，我们已经简要讨论过这个工具。\texttt{go doc} 命令会给出 \texttt{go} 命令能在子目录中找到的包或符号的文档。它有一个小限制：\texttt{go doc} 不会到互联网上查找，因为它是本地工具。这意味着你需要在一个项目中工作（其中有 \texttt{go.mod} 文件），而且该项目的依赖项已经下载。有些 IDE 会悄悄完成这件事，但你始终可以手动运行 \texttt{go mod download}。在这里，我们运行以下命令，分别取得 \texttt{pocketlog} 包、\texttt{pocketlog} 包中的 \texttt{New} 函数等内容的文档：

\begin{shellcode}[]
> go doc path/to/repo/pocketlog
> go doc path/to/repo/pocketlog.Logger
> go doc path/to/repo/pocketlog.New
> go doc path/to/repo/pocketlog.Logger.Debugf
\end{shellcode}
无论如何，文档都应始终以注释、示例或包头的形式成为交付内容的一部分。你可以在这里进一步了解：\url{https://go.dev/doc/comment}。既然已经说明了如何使用这个库，是时候让它真正可用了！

\section{实现导出的方法}

在上一节中，我们把 API 发布给 Susan，并编写了一些失败的测试。为了满足预期，下一步是决定日志记录器究竟在哪里执行任务，最后真正写下日志。我们会问这样的问题：它是否应该总是写入标准输出，绝不写入标准错误
% SOURCE_PDF_PAGE: 068
输出？它如何把消息发到票据打印机、写入文件，或通过网络发送给其他聚合机制？我们是否应该为每一种可能的用例实现这项功能？

我们希望这个包能在任何场景中使用，因此会把自定义写入器的实现留给用户自行决定。但 Susan 想要一个默认实现，直接把内容输出到控制台（标准输出），这样她就能先专注于业务逻辑，之后再改进日志记录和监控（我们可能并不认同这里的优先级，但产品团队坚持如此）。那就用标准输出吧。很快我们还会改进它。

\subsection{默认实现}

第一个实现是容易的部分（见代码清单 4.6）。在用下面的答案破坏你的解题乐趣之前，先想想自己会如何编写 \texttt{Debugf()} 方法。请记住：只有当阈值级别为 \texttt{Debug} 或更低时，\texttt{Debugf()} 才应记录日志。

\begin{gocode}[title={代码清单 4.6 logger.go：控制台的默认 Debugf 实现}]
// Debug 在日志级别为 debug 或更高时格式化并打印消息。
func (l *Logger) Debugf(format string, args ...any) {
    if l.threshold > LevelDebug { // (1)
        return
    }

    _, _ = fmt.Printf(format+"\n", args...) // (2)
}
\end{gocode}

\begin{codeannotations}
  \item \texttt{Debug} 会记录任何哪怕只有一点点意思的内容。
  \item 我们使用 \texttt{Printf}，它会把参数格式化成字符串。
\end{codeannotations}

调用 \texttt{Debugf} 函数时，用户期望：只要日志记录器的级别允许，消息就会被打印。因此，这个函数首先要确认是否应记录消息。我们声明的级别枚举让
% SOURCE_PDF_PAGE: 069
两个级别可以相互比较，因为底层的 \texttt{Level} 类型是整数。

如果选择在 \texttt{if} 内记录日志并反转条件，这个方法本可以只有三行；但你始终应该让正常路径保持不缩进。把错误处理和提前退出放在 \texttt{if} 块中，让真正的业务逻辑尽量靠左。这样做会显著提高代码可读性，也让扩展容易得多。

确认需要处理这条消息后，我们就记录它。目前先使用 \texttt{fmt.Printf} 函数。整个库看起来可能只是对这次 \texttt{Printf} 调用做了一层啰嗦的包装，但请放心，它并没有表面上那么简单。

\begin{infobox}{支线任务 4.3}
实现 \texttt{Info}、\texttt{Infof}、\texttt{Error} 和 \texttt{Errorf} 各级别的方法。此外，把你选择添加的所有其他级别也实现出来。
\end{infobox}

现在，先前的测试应该通过了。让我们悄悄推迟对其他方法的测试；我们想编写灵活性更高的 \texttt{TestXxx} 方法，因此需要写入非标准输出。

\subsection{使用接口}

把字节写到不同位置，是所有计算机程序中极其常见的用例。无论是把 JSON 写入 HTTP 输出、把 0 和 1 写入网络路由器、把比特写入数字端口来点亮一盏灯、把编码后的像素写入打印机，还是把字母写到控制台，我们所做的一切都必须把结果输出到某处才有用。Go 为最常见的用途提供了一组标准接口，例如写入或
% SOURCE_PDF_PAGE: 070
读取字节。这意味着，任何编写字节读写代码的人都可以使用这些现有接口。

\subsubsection*{io.Writer}

在标准库中最常被提及的接口里，\texttt{io} 包包含两个著名接口：\texttt{io.Writer}，用于写入任意目标；以及 \texttt{io.Reader}，用于从任意来源读取（例如字节数组、文件或 JSON 流）。下面的代码清单展示了这些接口的声明。

\begin{gocode}[title={代码清单 4.7 io 接口}]
type Reader interface {
    Read(p []byte) (n int, err error)
}

type Writer interface {
    Write(p []byte) (n int, err error)
}
\end{gocode}

我们希望日志记录器的用户定义目标位置。可以向用户索取一个 \texttt{io.Writer}，然后直接写入其中。用户负责提供自己选择的实现。

\begin{tipbox}{提示}
Go 与最著名的面向对象语言（例如 Java 和 C++）之间的一项主要区别是：Go 中的接口是隐式的。要实现一个接口，只需把接口定义的方法添加到对象上，编译器就会识别它。
\end{tipbox}

现在就可以给结构体添加输出，并把标准 \texttt{Writer} 添加到 \texttt{New()} 构建函数中。代码清单如下。

% SOURCE_PDF_PAGE: 071
\begin{gocode}[title={代码清单 4.8 logger.go：给结构体添加输出字段}]
// Logger 用于记录信息。
type Logger struct {
    threshold Level
    output    io.Writer
}

// New 返回一个日志记录器，可按所需阈值开始记录日志。
// 默认输出为 Stdout。
func New(threshold Level, output io.Writer) *Logger {
    return &Logger{
        threshold: threshold,
        output:    output,
    }
}
\end{gocode}

注意这段文档略有增强。也可以把默认值设为 \texttt{os.Stderr}，它表示默认错误输出。

最后，每个方法都需要使用这个新字段（见代码清单 4.9）。还要确保那些没有遵循我们的建议、未用 \texttt{New} 创建 \texttt{Logger} 的用户不会遇到空指针引用异常。

\begin{gocode}[title={代码清单 4.9 logger.go：给结构体添加输出字段}]
// Debug 在日志级别为 debug 或更高时格式化并打印消息。
func (l Logger) Debugf(format string, args ...any) {
    // 确保可以安全地写入输出
    if l.output == nil {
        l.output = os.Stdout
    }
    if l.threshold <= LevelDebug {
        _, _ = fmt.Fprintf(l.output, format, args...)
    }
}
\end{gocode}

在 Go 中，下划线符号表示 void（无值）。换句话说，把一个值赋给下划线只会丢弃结果。这样做有什么意义？Go 喜欢明确表达。这里，我们明确告诉下一位开发人员，也包括未来的自己：“我知道这个函数会返回一些值，但我不需要它们。”

这里，函数返回一个 \texttt{int}（写入的字符数），有时还返回一个错误。此刻我们不想处理
% SOURCE_PDF_PAGE: 072
这个错误，所以明确地忽略它。

\subsection{重构}

在实现 \texttt{Info} 和 \texttt{Error} 方法时，你可能已经注意到，我们调用的是同一个写入函数 \texttt{fmt.\allowbreak Fprintf}。你还可能注意到，与同属 \texttt{fmt} 包的 \texttt{fmt.Fprintf} 不同，\texttt{fmt.Fprintln} 会在字符串末尾追加一个行尾字符。从某种意义上说，\texttt{printf} 函数允许你进行非常精细的格式控制，而 \texttt{println} 函数不会让你完全按照自己的喜好格式化一切：它无法在数字或字符串前补齐空格，也无法使用数字的十六进制表示，等等。

换行可以保证日志消息在控制台中容易区分。既然我们想导出 \texttt{Printf} 的整套工具箱，就必须在写入消息时显式添加 \texttt{\textbackslash n}。

在这里，这个换行需要添加三次，也就是在三个函数中各加一次。而且每当我们想修改日志消息（例如添加日志级别）时，都必须把相同的几行写三遍。如果还要实现 \texttt{Warn} 或 \texttt{Fatal}，重复次数会更多。这（大声地）呼唤一次小型重构；我们不想维护三份甚至更多份相同的代码。让我们把所有这些行归在一起，这样每次调整日志记录器时只需改一行。

在 \texttt{Logger} 上创建一个 \texttt{logf()} 方法，如代码清单 4.10 所示。暂时让该方法接受与 \texttt{Debug}、\texttt{Info} 和 \texttt{Error} 相同的参数。这三个方法将调用 \texttt{logf}；现在由 \texttt{logf} 独自负责格式化和打印。另外三个导出方法则各自只负责自己的日志
% SOURCE_PDF_PAGE: 073
级别，除此之外什么也不管。没有任何充分理由导出 \texttt{logf}。

\begin{gocode}[title={代码清单 4.10 logger.go：重构日志记录方法}]
// Debugf 在日志级别为 debug 或更高时格式化并打印消息。
func (l *Logger) Debugf(format string, args ...any) {
    if l.threshold > LevelDebug {
        return
    }

    l.logf(format, args...)
}

// logf 把消息打印到输出。
// 如有修饰内容，请在这里添加。(1)
func (l *Logger) logf(format string, args ...any) {
    _, _ = fmt.Fprintf(l.output, format+"\n", args...)
}
\end{gocode}

\begin{codeannotations}
  \item 未导出方法上的注释通常是写给下一位维护人员的（包括你自己）。
\end{codeannotations}

运行测试。现在测试应该会失败，因为我们还没有用 \texttt{New} 函数的新参数 \texttt{os.Stdout} 更新它们。完成更新后，你就可以提交，并告诉同事代码已经可以使用。然而 Susan 告诉你：她的日志应该写入一个特定文件，而不是标准输出，因为她已经在使用标准输出了。\texttt{Logger} 自己能做到吗？

\section{函数式选项模式}

默认值是开发理论家们激烈争论的主题。有些人喜欢使用默认值，因为它让一切都易于发现，因此更容易起步。另一些人讨厌默认值，因为它不会迫使你弄清楚自己在做什么。我们认为，只要适度使用，默认值就是好东西。

比如，当我们开始开发一项新的 Web 服务时，希望专注于业务逻辑。我们要确保自己的
% SOURCE_PDF_PAGE: 074
逻辑能在本地运行，然后才让服务达到生产就绪状态。为了降低认知负担，我们从日志记录器的默认版本开始，更好的日志记录器架构可以稍后再做。在编写部署代码之前，除了标准输出，我们什么也不需要。

但很快，我们就会把它部署到一家便宜的云服务商，以便推介原型并向全世界展示。此时，读取标准输出就没那么简单了。我们选用一个工具，比如某种聚合数据库，而它恰好发布了一个 Go 驱动。这个驱动的开发人员很聪明，在库中提供了一个实现 \texttt{io.Writer} 接口的结构体。

我们仍想保留这个写入标准输出的默认实现，但也希望提供写到其他位置的选项。实现这一点的一种常见方式，是使用函数式选项模式。

\subsection{创建配置}

在 \texttt{pocketlog} 包中新建一个 \texttt{options.go} 文件，用来编写配置函数。定义一种函数类型，它可传给 \texttt{New()} 函数，并会依次逐个应用：

\begin{gocode}[]
// Option 为日志记录器定义一个函数式选项。
type Option func(*Logger)
\end{gocode}

这个函数接收指向日志记录器的指针，因此函数体可以直接修改指针所指的值。在我们的例子中，就是把默认输出改为用户给出的任何内容。

% SOURCE_PDF_PAGE: 075
\begin{gocode}[title={代码清单 4.11 options.go：更改输出的可选函数}]
// WithOutput 返回一个用于设置日志输出的配置函数。
func WithOutput(output io.Writer) Option {
    return func(lgr *Logger) {
        lgr.output = output
    }
}
\end{gocode}

这种函数可以作为可变参数传给 \texttt{New()} 函数，也就是由零个或多个同类型参数组成的列表。

\begin{gocode}[title={代码清单 4.12 logger.go：应用选项}]
// New 返回一个日志记录器，可按所需阈值开始记录日志。
// 给它一组配置函数，即可随心调整。
// 默认输出为 Stdout。
func New(threshold Level, opts ...Option) *Logger {
    lgr := &Logger{threshold: threshold, output: os.Stdout}

    for _, configFunc := range opts {
        configFunc(lgr)
    }

    return lgr
}
\end{gocode}

下次想给日志记录器添加选项时（例如日期格式化器），只需创建一个新的 \texttt{Option}，一切就绪。这里很重要的一点是：添加配置函数很容易，它让用户无需改变库的 API 就能设置特定行为。我们的 \texttt{New} 函数可以接受用户所需数量的配置函数，这些函数来自我们在本包中实现的列表。

\subsubsection*{使用示例}

Susan 想知道如何使用你的库。虽然已经有文档文件，但人与人之间的交流总是高效得多。于是，你编写一个小示例发给她。

在库之外，运行 \texttt{go mod init} 初始化一个新模块，并创建 \texttt{main.go} 文件。像上一
% SOURCE_PDF_PAGE: 076
章一样定义一个 \texttt{main()} 函数。在这个函数中实例化一个新的日志记录器，再调用几个方法来展示你的成果。

\begin{gocode}[title={代码清单 4.13 main.go：使用示例}]
package main

import (
    "os"
    "time"

    "learngo-pockets/logger/pocketlog"
)

func main() {
    lgr := pocketlog.New(pocketlog.LevelInfo, pocketlog.WithOutput(os.
        Stdout))

    lgr.Infof("A little copying is better than a little dependency.")
    lgr.Errorf("Errors are values. Documentation is for %s.", "users")
    lgr.Debugf("Make the zero (%d) value useful.", 0)

    lgr.Infof("Hallo, %d %v", 2022, time.Now())
}
\end{gocode}

\subsection{测试我们的示例}

我们已经在使用日志记录器，但它还没有经过完整测试。太不专业了！Susan 可以使用我们的库，但我们可不想让她带着可能存在的缺陷回来。接口的神奇之处意味着：我们可以编写一个实现 \texttt{io.Writer} 的测试辅助工具，并把它交给受测的 \texttt{Logger}。

\subsubsection*{测试辅助工具的实现}

在 \texttt{logger\_test.go} 末尾编写一个新的 \texttt{testWriter} 结构体，如代码清单 4.14 所示。让它实现 \texttt{io.Writer}；但它不把内容写到某个目标，而是验证输出字符串。例如，可以在结构体中保留一个字段，把输出拼接到其中，再拿它与预期结果比较。

% SOURCE_PDF_PAGE: 077
\begin{gocode}[title={代码清单 4.14 \texttt{logger\_test.go}：测试辅助工具的实现}]
// testWriter 是一个实现 io.Writer 的结构体。
// 我们用它验证能否写入特定输出。
type testWriter struct {
    contents string
}

// Write 实现 io.Writer 接口。
func (tw *testWriter) Write(p []byte) (n int, err error) { // (1)
    tw.contents = tw.contents + string(p) // (2)
    return len(p), nil
}
\end{gocode}

\begin{codeannotations}
  \item 接收器是一个指针，因此我们可以修改它的 \texttt{contents} 字段。
  \item 把每次调用的内容拼接到内存中的日志记录器里。
\end{codeannotations}

在测试前面的部分，可以把这个结构体传给函数式选项。测试结束时，就可以检查写入器的内容是否符合预期。

实际使用中，可以用 \texttt{strings.Builder} 或 \texttt{bytes.Buffer} 代替 \texttt{testWriter}。现在，即使需要的接口不属于标准库，你也知道怎样制作模拟对象（mock）了。

\subsubsection*{更新测试}

现在不再被迫检查标准输出，我们可以编写一个 \texttt{TestXxx} 函数，依次把所有日志记录方法一起测试（见代码清单 4.15）。每个必需的日志级别都可以有一个测试用例；我们可以检查输出是否不同，并确认 \texttt{Debugf()} 调用大多会被忽略。

% SOURCE_PDF_PAGE: 078
\begin{gocode}[title={代码清单 4.15 \texttt{logger\_test.fo}：测试函数}]
const ( // (1)
    debugMessage = "Why write I still all one, ever the same,"
    infoMessage = "And keep invention in a noted weed,"
    errorMessage = "That every word doth almost tell my name,"
)

func TestLogger_DebugfInfofErrorf(t *testing.T) {
    type testCase struct {
        level    pocketlog.Level
        expected string
    }

    tt := map[string]testCase{
        "debug": {
            level:    pocketlog.LevelDebug,
            expected: debugMessage + "\n" +
                      infoMessage + "\n" +
                      errorMessage + "\n",
        },
        "info": {...}, // (2)
        "error": {...},
    }

    for name, tc := range tt {
        t.Run(name, func(t *testing.T) {
            tw := &testWriter{}

            testedLogger := pocketlog.New(tc.level, pocketlog.WithOutput(tw))

            testedLogger.Debugf(debugMessage)
            testedLogger.Infof(infoMessage)
            testedLogger.Errorf(errorMessage)

            if tw.contents != tc.expected {
                t.Errorf("invalid contents, expected %q, got %q",
                    tc.expected, tw.contents) // (3)
            }
        })
    }
}
\end{gocode}

\begin{codeannotations}
  \item 定义测试字符串；请使用你最喜欢的十四行诗。
  \item 轮到你编写这两个测试用例了！
  \item 注意使用 \texttt{\%q}，它有助于检查多余空格。
\end{codeannotations}

% SOURCE_PDF_PAGE: 079
\begin{infobox}{支线任务 4.4}
按我们这里的写法，这个测试只会按一种顺序测试函数调用：先 \texttt{Debugf}，再 \texttt{Infof}，然后是 \texttt{Errorf}。如果我们决定添加一个缓冲区，并且只在 \texttt{Errorf()} 方法中才把所有内容写出去，会怎么样？当前测试发现不了这个问题，而 \texttt{Debug} 和 \texttt{Info} 消息可能会一直卡在缓冲区里。
\end{infobox}

你的日志记录器已经准备就绪，功能完备、有文档且经过测试。公司的其他人开始使用它，而你还在不断构想新功能。让我们探索其中一些，看看它们会把我们带到哪里。

\section{其他功能}

这个工具有无穷无尽的优化可能。和往常一样，唯一的限制是我们愿意投入多少时间。例如，这个库并非线程安全；当多个 goroutine 在没有任何保护的情况下使用同一个 \texttt{Writer} 时，结果可能出乎意料。我们将在后面的章节中探索解决方案。与此同时，先看看可以给 \texttt{Logger} 添加的一些有趣增强功能。

\subsection{记录日志级别}

你的服务在本地运行，采用尽可能低的日志级别。只要看着控制台，你就知道发生的一切；但现在，你希望以红色显示错误。你在日志跟踪命令中加入了一条老式的 \texttt{awk} 命令，可你怎么知道该给什么内容上色？
